% Options for packages loaded elsewhere
\PassOptionsToPackage{unicode}{hyperref}
\PassOptionsToPackage{hyphens}{url}
%
\documentclass[
]{article}
\usepackage{lmodern}
\usepackage{amssymb,amsmath}
\usepackage{ifxetex,ifluatex}
\ifnum 0\ifxetex 1\fi\ifluatex 1\fi=0 % if pdftex
  \usepackage[T1]{fontenc}
  \usepackage[utf8]{inputenc}
  \usepackage{textcomp} % provide euro and other symbols
\else % if luatex or xetex
  \usepackage{unicode-math}
  \defaultfontfeatures{Scale=MatchLowercase}
  \defaultfontfeatures[\rmfamily]{Ligatures=TeX,Scale=1}
\fi
% Use upquote if available, for straight quotes in verbatim environments
\IfFileExists{upquote.sty}{\usepackage{upquote}}{}
\IfFileExists{microtype.sty}{% use microtype if available
  \usepackage[]{microtype}
  \UseMicrotypeSet[protrusion]{basicmath} % disable protrusion for tt fonts
}{}
\makeatletter
\@ifundefined{KOMAClassName}{% if non-KOMA class
  \IfFileExists{parskip.sty}{%
    \usepackage{parskip}
  }{% else
    \setlength{\parindent}{0pt}
    \setlength{\parskip}{6pt plus 2pt minus 1pt}}
}{% if KOMA class
  \KOMAoptions{parskip=half}}
\makeatother
\usepackage{xcolor}
\IfFileExists{xurl.sty}{\usepackage{xurl}}{} % add URL line breaks if available
\IfFileExists{bookmark.sty}{\usepackage{bookmark}}{\usepackage{hyperref}}
\hypersetup{
  pdftitle={Using CACE for a Chipalooza CMOS5L Schematic-Review Submission},
  pdfauthor={Compiled for Christoph Maier},
  hidelinks,
  pdfcreator={LaTeX via pandoc}}
\urlstyle{same} % disable monospaced font for URLs
\usepackage[margin=1in]{geometry}
\usepackage{color}
\usepackage{fancyvrb}
\newcommand{\VerbBar}{|}
\newcommand{\VERB}{\Verb[commandchars=\\\{\}]}
\DefineVerbatimEnvironment{Highlighting}{Verbatim}{commandchars=\\\{\}}
% Add ',fontsize=\small' for more characters per line
\newenvironment{Shaded}{}{}
\newcommand{\AlertTok}[1]{\textcolor[rgb]{1.00,0.00,0.00}{\textbf{#1}}}
\newcommand{\AnnotationTok}[1]{\textcolor[rgb]{0.38,0.63,0.69}{\textbf{\textit{#1}}}}
\newcommand{\AttributeTok}[1]{\textcolor[rgb]{0.49,0.56,0.16}{#1}}
\newcommand{\BaseNTok}[1]{\textcolor[rgb]{0.25,0.63,0.44}{#1}}
\newcommand{\BuiltInTok}[1]{#1}
\newcommand{\CharTok}[1]{\textcolor[rgb]{0.25,0.44,0.63}{#1}}
\newcommand{\CommentTok}[1]{\textcolor[rgb]{0.38,0.63,0.69}{\textit{#1}}}
\newcommand{\CommentVarTok}[1]{\textcolor[rgb]{0.38,0.63,0.69}{\textbf{\textit{#1}}}}
\newcommand{\ConstantTok}[1]{\textcolor[rgb]{0.53,0.00,0.00}{#1}}
\newcommand{\ControlFlowTok}[1]{\textcolor[rgb]{0.00,0.44,0.13}{\textbf{#1}}}
\newcommand{\DataTypeTok}[1]{\textcolor[rgb]{0.56,0.13,0.00}{#1}}
\newcommand{\DecValTok}[1]{\textcolor[rgb]{0.25,0.63,0.44}{#1}}
\newcommand{\DocumentationTok}[1]{\textcolor[rgb]{0.73,0.13,0.13}{\textit{#1}}}
\newcommand{\ErrorTok}[1]{\textcolor[rgb]{1.00,0.00,0.00}{\textbf{#1}}}
\newcommand{\ExtensionTok}[1]{#1}
\newcommand{\FloatTok}[1]{\textcolor[rgb]{0.25,0.63,0.44}{#1}}
\newcommand{\FunctionTok}[1]{\textcolor[rgb]{0.02,0.16,0.49}{#1}}
\newcommand{\ImportTok}[1]{#1}
\newcommand{\InformationTok}[1]{\textcolor[rgb]{0.38,0.63,0.69}{\textbf{\textit{#1}}}}
\newcommand{\KeywordTok}[1]{\textcolor[rgb]{0.00,0.44,0.13}{\textbf{#1}}}
\newcommand{\NormalTok}[1]{#1}
\newcommand{\OperatorTok}[1]{\textcolor[rgb]{0.40,0.40,0.40}{#1}}
\newcommand{\OtherTok}[1]{\textcolor[rgb]{0.00,0.44,0.13}{#1}}
\newcommand{\PreprocessorTok}[1]{\textcolor[rgb]{0.74,0.48,0.00}{#1}}
\newcommand{\RegionMarkerTok}[1]{#1}
\newcommand{\SpecialCharTok}[1]{\textcolor[rgb]{0.25,0.44,0.63}{#1}}
\newcommand{\SpecialStringTok}[1]{\textcolor[rgb]{0.73,0.40,0.53}{#1}}
\newcommand{\StringTok}[1]{\textcolor[rgb]{0.25,0.44,0.63}{#1}}
\newcommand{\VariableTok}[1]{\textcolor[rgb]{0.10,0.09,0.49}{#1}}
\newcommand{\VerbatimStringTok}[1]{\textcolor[rgb]{0.25,0.44,0.63}{#1}}
\newcommand{\WarningTok}[1]{\textcolor[rgb]{0.38,0.63,0.69}{\textbf{\textit{#1}}}}
\usepackage{longtable,booktabs}
% Correct order of tables after \paragraph or \subparagraph
\usepackage{etoolbox}
\makeatletter
\patchcmd\longtable{\par}{\if@noskipsec\mbox{}\fi\par}{}{}
\makeatother
% Allow footnotes in longtable head/foot
\IfFileExists{footnotehyper.sty}{\usepackage{footnotehyper}}{\usepackage{footnote}}
\makesavenoteenv{longtable}
\setlength{\emergencystretch}{3em} % prevent overfull lines
\providecommand{\tightlist}{%
  \setlength{\itemsep}{0pt}\setlength{\parskip}{0pt}}
\setcounter{secnumdepth}{-\maxdimen} % remove section numbering

\title{Using CACE for a Chipalooza CMOS5L Schematic-Review Submission}
\usepackage{etoolbox}
\makeatletter
\providecommand{\subtitle}[1]{% add subtitle to \maketitle
  \apptocmd{\@title}{\par {\large #1 \par}}{}{}
}
\makeatother
\subtitle{An excerpt-and-walkthrough guide, adapted from the JKU
\emph{Analog (Integrated) Circuit Design} course}
\author{Compiled for Christoph Maier}
\date{22 August 2026}

\begin{document}
\maketitle

{
\setcounter{tocdepth}{3}
\tableofcontents
}
\begin{quote}
\textbf{Math note.} This document uses \texttt{\$...\$} for inline math
and \texttt{\$\$...\$\$} for display math, which renders on GitHub.
GitLab's default Markdown renderer expects
\texttt{\$\textasciigrave{}...\textasciigrave{}\$} for inline math
instead of \texttt{\$...\$}; if you view this file on GitLab and the
inline equations look wrong, that's why --- the display equations
(\texttt{\$\$...\$\$}) will still render correctly on both.
\end{quote}

\begin{quote}
\textbf{Two deadlines, both 31 August 2026, two different deliverable
types:}

\begin{longtable}[]{@{}lll@{}}
\toprule
\begin{minipage}[b]{0.30\columnwidth}\raggedright
Track\strut
\end{minipage} & \begin{minipage}[b]{0.30\columnwidth}\raggedright
Milestone due 31 Aug 2026\strut
\end{minipage} & \begin{minipage}[b]{0.30\columnwidth}\raggedright
What's actually required\strut
\end{minipage}\tabularnewline
\midrule
\endhead
\begin{minipage}[t]{0.30\columnwidth}\raggedright
\textbf{SG13CMOS5L}\strut
\end{minipage} & \begin{minipage}[t]{0.30\columnwidth}\raggedright
\textbf{Schematic review}\strut
\end{minipage} & \begin{minipage}[t]{0.30\columnwidth}\raggedright
A committed CACE run
(\texttt{docs/\textless{}name\textgreater{}\_schematic.md}) against your
xschem schematic --- §§1--8 and §13 of this document are your critical
path. This is the deliverable this walkthrough is built around.\strut
\end{minipage}\tabularnewline
\begin{minipage}[t]{0.30\columnwidth}\raggedright
\textbf{GF180MCU (``180u'')}\strut
\end{minipage} & \begin{minipage}[t]{0.30\columnwidth}\raggedright
\textbf{Proposal}\strut
\end{minipage} & \begin{minipage}[t]{0.30\columnwidth}\raggedright
A written proposal document (same format as your
\texttt{s2d\_d2s\_pinbuffers.md}) --- no CACE run required yet. §9.2
(GF180MCU adaptation) is \emph{get-ahead} reference for after the
proposal is accepted, not something you need to finish by the 31st for
this track.\strut
\end{minipage}\tabularnewline
\bottomrule
\end{longtable}

Don't let GF180MCU CACE/PDK-porting work compete for time against the
CMOS5L schematic-review evidence this week --- the GF180MCU track only
needs prose and a circuit description right now.
\end{quote}

\hypertarget{contents}{%
\subsection{Contents}\label{contents}}

\begin{enumerate}
\def\labelenumi{\arabic{enumi}.}
\tightlist
\item
  \protect\hyperlink{1-what-cace-is-and-why-it-matters-for-schematic-review}{What
  CACE Is, and Why It Matters for Schematic Review}
\item
  \protect\hyperlink{2-cace-project-layout}{CACE Project Layout}
\item
  \protect\hyperlink{3-anatomy-of-the-yaml-datasheet}{Anatomy of the
  YAML Datasheet}
\item
  \protect\hyperlink{4-testbench-templates-and-token-substitution}{Testbench
  Templates and Token Substitution}
\item
  \protect\hyperlink{5-running-cace-and-reading-the-report}{Running CACE
  and Reading the Report}
\item
  \protect\hyperlink{6-pvt-corner-sweeps}{PVT Corner Sweeps}
\item
  \protect\hyperlink{7-monte-carlo-and-mismatch-analysis}{Monte Carlo
  and Mismatch Analysis}
\item
  \protect\hyperlink{8-post-layout-characterization}{Post-Layout
  Characterization}
\item
  \protect\hyperlink{9-per-pdk-adaptation-sg13g2cmos5l-gf180mcu-sky130}{Per-PDK
  Adaptation: SG13G2/CMOS5L, GF180MCU, SKY130}
\item
  \protect\hyperlink{10-model-accuracy-caveats}{Model-Accuracy Caveats}
\item
  \protect\hyperlink{11-mapping-cace-output-to-chipalooza-schematic-review-expectations}{Mapping
  CACE Output to Chipalooza Schematic-Review Expectations}
\item
  \protect\hyperlink{12-appendix-command-cheat-sheet-and-blank-datasheet-skeleton}{Appendix:
  Command Cheat-Sheet and Blank Datasheet Skeleton}
\item
  \protect\hyperlink{13-the-chipalooza-2-analog-harness-and-what-it-means-for-your-cace-setup}{The
  Chipalooza \#2 Analog Harness, and What It Means for Your CACE Setup}
\end{enumerate}

\begin{center}\rule{0.5\linewidth}{0.5pt}\end{center}

\hypertarget{what-cace-is-and-why-it-matters-for-schematic-review}{%
\subsection{1. What CACE Is, and Why It Matters for Schematic
Review}\label{what-cace-is-and-why-it-matters-for-schematic-review}}

CACE (\href{https://github.com/efabless/cace}{Circuit Automatic
Characterization Engine}, maintained by efabless) is a Python-based
simulation runner for analog and mixed-signal circuits. You write one
specification file --- the \textbf{datasheet}, in YAML --- that lists
every pin, every operating condition (supply voltage, temperature,
process corner, load, \ldots), and every electrical parameter you're
claiming for the circuit, together with pass/fail limits. CACE then
generates the SPICE netlists, runs ngspice for every condition
combination, extracts the results, checks them against your limits, and
writes a Markdown/HTML report with plots.

Two things make this directly relevant to a Chipalooza schematic-review
deadline rather than a ``nice to have'':

\begin{itemize}
\item
  \textbf{The schematic-review artifact \emph{is} a CACE report.}
  Efabless's own analog IP convention (see
  \protect\hyperlink{2-cace-project-layout}{§2}) expects a
  \texttt{docs/\textless{}name\textgreater{}\_schematic.md} file
  generated straight from the CACE run against your xschem schematic. A
  reviewer opening your repo expects to find your claimed specs, the
  conditions you tested them under, and a pass/fail table --- not just a
  schematic picture and a promise.
\item
  \textbf{It forces you to write the spec before/while you simulate},
  rather than screenshotting a few waveforms after the fact. As your
  course material puts it (quoting the \emph{Analog (Integrated) Circuit
  Design} textbook, §6.5):

  \begin{quote}
  ``Luckily, there are open-source versions of simulation runners
  available, and we will use
  \href{https://github.com/efabless/cace}{CACE} in this lecture. CACE is
  written in Python and allows one to set up a datasheet in
  \href{https://yaml.org}{YAML} that defines the simulation problem and
  the performance parameters to evaluate against specified limits.'' ---
  H. Pretl et al., \emph{Analog (Integrated) Circuit Design}, §``PVT
  Simulation and Design Corners''
  \end{quote}
\end{itemize}

CACE's own documentation
(\href{https://cace.readthedocs.io/en/latest/}{cace.readthedocs.io})
frames the same point as a workflow benefit: \emph{``CACE encourages
good analog design practices, fostering trust in open source analog
design.''} For a multi-team, multi-block shuttle like Chipalooza, that
trust is the entire point of the schematic review --- the reviewer
cannot re-derive your circuit's operating point by eye across 60+
submissions, but they can read a CACE table.

\begin{center}\rule{0.5\linewidth}{0.5pt}\end{center}

\hypertarget{cace-project-layout}{%
\subsection{2. CACE Project Layout}\label{cace-project-layout}}

Efabless publishes a template repository,
\href{https://github.com/efabless/sky130_ef_ip__template}{\texttt{sky130\_ef\_ip\_\_template}},
that fixes the directory convention every Chipalooza-style analog IP
block is expected to follow. A real Chipalooza entry ---
\href{https://github.com/efabless/sky130_sw_ip__por}{\texttt{sky130\_sw\_ip\_\_por}},
a bandgap-referenced Power-on-Reset block --- follows it too. Your
course's own \texttt{cace/} folder (for the 5T-OTA and improved-OTA
examples) is the same convention, just embedded inside a textbook
chapter instead of standing alone as a repo.

\begin{verbatim}
<pdk_family>_<handle>_ip__<project>/
├── cace/          — CACE datasheet (YAML) + testbench templates
├── docs/          — CACE-generated reports:
│                    <name>_schematic.md, _layout.md, _pex.md, _rcx.md
├── xschem/        — <name>.sch / <name>.sym (top level), <name>_tb.sch (testbench),
│                    <name>/ (internal subcircuits), xschemrc
├── netlist/       — SPICE netlists CACE generates (subfolders: schematic/layout/pex/rcx)
├── mag/           — Magic layout (optional; delete if unused)
├── gds/           — GDSII
├── ip/            — dependency IP blocks, added as git submodules
├── runs/          — CACE run output (gitignored)
├── .github/workflows/cace.yml   — CI: re-runs CACE automatically on push
└── README.md      — links to docs/, one line per characterization stage
\end{verbatim}

Repository naming convention:
\textbf{\texttt{\textless{}pdk\_family\textgreater{}\_\textless{}handle\textgreater{}\_ip\_\_\textless{}project\textgreater{}}}
(double underscore after \texttt{ip}). The
\texttt{\textless{}handle\textgreater{}} is a short unique identifier
--- your initials work --- to avoid collisions between different
designers' blocks of the same name. For this project that would be
something like \texttt{ihp13g2\_cm\_ip\_\_s2d\_d2s\_pinbuffers}.

Two structural points worth internalizing before you build your own
tree:

\begin{itemize}
\tightlist
\item
  \textbf{\texttt{docs/} is stage-tagged, not just one report.} CACE can
  characterize the \emph{same} datasheet against four different netlist
  sources --- \texttt{schematic}, \texttt{layout} (LVS), \texttt{pex}
  (C-parasitic extracted), \texttt{rcx} (full R-C extracted) ---
  selected with the \texttt{-s/-\/-source} flag (see
  \protect\hyperlink{8-post-layout-characterization}{§8}). The
  convention is one report file per stage:
  \texttt{\textless{}name\textgreater{}\_schematic.md} is what a
  \emph{schematic} reviewer wants; \texttt{\_pex.md}/\texttt{\_rcx.md}
  come later once you have layout.
\item
  \textbf{\texttt{ip/} is how blocks depend on each other}, added as git
  submodules
  (\texttt{git\ submodule\ add\ ../\textless{}ip\_block\textgreater{}.git\ ip/\textless{}ip\_block\textgreater{}}),
  with \texttt{xschemrc} sourcing each dependency's \texttt{xschemrc}
  recursively. If your gmC filter block reuses an Oguey--Aebischer bias
  cell as a separate IP, this is the mechanism, not a copy-paste.
\end{itemize}

\begin{center}\rule{0.5\linewidth}{0.5pt}\end{center}

\hypertarget{anatomy-of-the-yaml-datasheet}{%
\subsection{3. Anatomy of the YAML
Datasheet}\label{anatomy-of-the-yaml-datasheet}}

The datasheet format (current version \texttt{5.2}, per the
\href{https://cace.readthedocs.io/en/latest/reference/datasheet_format.html}{CACE
reference manual}) has six top-level sections. Below, each is explained
against the \textbf{real, working datasheet from your course} ---
\texttt{cace/voltage-buffer-ota.yaml}, characterizing the 5T-OTA voltage
buffer --- so this is an excerpt of a file that already runs, not a
hypothetical.

\hypertarget{metadata-and-authorship}{%
\subsubsection{3.1 Metadata and
authorship}\label{metadata-and-authorship}}

\begin{Shaded}
\begin{Highlighting}[]
\FunctionTok{name}\KeywordTok{:}\AttributeTok{           ota{-}5t}
\FunctionTok{description}\KeywordTok{:}\AttributeTok{    Simple voltage buffer for capacitive load realized with 5T{-}OTA}
\FunctionTok{PDK}\KeywordTok{:}\AttributeTok{            ihp{-}sg13g2}\CommentTok{ \# (for IIC{-}OSIC{-}Tools version 2025.01 or newer)}
\FunctionTok{cace\_format}\KeywordTok{:}\AttributeTok{    }\FloatTok{5.2}

\FunctionTok{authorship}\KeywordTok{:}
\AttributeTok{  }\FunctionTok{designer}\KeywordTok{:}\AttributeTok{         Harald Pretl}
\AttributeTok{  }\FunctionTok{company}\KeywordTok{:}\AttributeTok{          Johannes Kepler University}
\AttributeTok{  }\FunctionTok{creation\_date}\KeywordTok{:}\AttributeTok{    August 25, 2024}
\AttributeTok{  }\FunctionTok{license}\KeywordTok{:}\AttributeTok{          Apache 2.0}
\end{Highlighting}
\end{Shaded}

\texttt{name} must match the cell name; \texttt{PDK} is looked up
against your installed PDK (no spaces). This is also where you'd write
\texttt{ihp-sg13cmos5l} once that PDK identifier stabilizes (see
\protect\hyperlink{9-per-pdk-adaptation-sg13g2cmos5l-gf180mcu-sky130}{§9}).

\hypertarget{paths}{%
\subsubsection{3.2 Paths}\label{paths}}

\begin{Shaded}
\begin{Highlighting}[]
\FunctionTok{paths}\KeywordTok{:}
\AttributeTok{  }\FunctionTok{root}\KeywordTok{:}\AttributeTok{             ..}
\AttributeTok{  }\FunctionTok{schematic}\KeywordTok{:}\AttributeTok{        xschem}
\AttributeTok{  }\FunctionTok{netlist}\KeywordTok{:}\AttributeTok{          cace/netlist}
\AttributeTok{  }\FunctionTok{documentation}\KeywordTok{:}\AttributeTok{    cace/\_docs}
\AttributeTok{  }\FunctionTok{runs}\KeywordTok{:}\AttributeTok{             \_runs}
\end{Highlighting}
\end{Shaded}

\texttt{root} anchors everything else relative to the \texttt{cace/}
folder's parent --- normally your project root. \texttt{netlist} is
where CACE writes generated SPICE, subdivided automatically by netlist
source (\texttt{schematic/}, \texttt{layout/}, \texttt{pex/},
\texttt{rcx/}). \texttt{documentation} is where the Markdown/HTML report
lands --- this is your
\texttt{docs/\textless{}name\textgreater{}\_schematic.md} target.

\hypertarget{pins}{%
\subsubsection{3.3 Pins}\label{pins}}

\begin{Shaded}
\begin{Highlighting}[]
\FunctionTok{pins}\KeywordTok{:}
\AttributeTok{  }\FunctionTok{vdd}\KeywordTok{:}
\AttributeTok{    }\FunctionTok{description}\KeywordTok{:}\AttributeTok{ Positive analog power supply}
\AttributeTok{    }\FunctionTok{type}\KeywordTok{:}\AttributeTok{ power}
\AttributeTok{    }\FunctionTok{direction}\KeywordTok{:}\AttributeTok{ inout}
\AttributeTok{    }\FunctionTok{Vmin}\KeywordTok{:}\AttributeTok{ }\FloatTok{1.45}
\AttributeTok{    }\FunctionTok{Vmax}\KeywordTok{:}\AttributeTok{ }\FloatTok{1.55}
\AttributeTok{  }\FunctionTok{vss}\KeywordTok{:}
\AttributeTok{    }\FunctionTok{description}\KeywordTok{:}\AttributeTok{ Analog ground}
\AttributeTok{    }\FunctionTok{type}\KeywordTok{:}\AttributeTok{ ground}
\AttributeTok{    }\FunctionTok{direction}\KeywordTok{:}\AttributeTok{ inout}
\AttributeTok{  }\FunctionTok{ibias\_20u}\KeywordTok{:}
\AttributeTok{    }\FunctionTok{description}\KeywordTok{:}\AttributeTok{ Bias current input 20uA nom.}
\AttributeTok{    }\FunctionTok{type}\KeywordTok{:}\AttributeTok{ signal}
\AttributeTok{    }\FunctionTok{direction}\KeywordTok{:}\AttributeTok{ input}
\AttributeTok{  }\FunctionTok{vinp}\KeywordTok{:}
\AttributeTok{    }\FunctionTok{description}\KeywordTok{:}\AttributeTok{ Voltage positive input}
\AttributeTok{    }\FunctionTok{type}\KeywordTok{:}\AttributeTok{ signal}
\AttributeTok{    }\FunctionTok{direction}\KeywordTok{:}\AttributeTok{ input}
\AttributeTok{  }\FunctionTok{vinn}\KeywordTok{:}
\AttributeTok{    }\FunctionTok{description}\KeywordTok{:}\AttributeTok{ Voltage negative input}
\AttributeTok{    }\FunctionTok{type}\KeywordTok{:}\AttributeTok{ signal}
\AttributeTok{    }\FunctionTok{direction}\KeywordTok{:}\AttributeTok{ input}
\AttributeTok{  }\FunctionTok{vout}\KeywordTok{:}
\AttributeTok{    }\FunctionTok{description}\KeywordTok{:}\AttributeTok{ Voltage output}
\AttributeTok{    }\FunctionTok{type}\KeywordTok{:}\AttributeTok{ signal}
\AttributeTok{    }\FunctionTok{direction}\KeywordTok{:}\AttributeTok{ output}
\end{Highlighting}
\end{Shaded}

Pin names here are checked against your schematic/netlist directly --- a
typo means CACE fails loudly rather than silently simulating garbage.
\texttt{type} is one of \texttt{digital}, \texttt{signal},
\texttt{power}, \texttt{ground}; \texttt{direction} is
\texttt{input}/\texttt{output}/\texttt{inout}.
\texttt{Vmin}/\texttt{Vmax}/\texttt{Imin}/\texttt{Imax} can reference
other pins by expression (e.g.~\texttt{vdd\ +\ 0.3}), which is how you'd
express, say, an I/O pin's absolute maximum relative to a rail rather
than as a hardcoded number.

\hypertarget{default-conditions}{%
\subsubsection{3.4 Default conditions}\label{default-conditions}}

\begin{Shaded}
\begin{Highlighting}[]
\FunctionTok{default\_conditions}\KeywordTok{:}
\AttributeTok{  }\FunctionTok{vdd}\KeywordTok{:}
\AttributeTok{    }\FunctionTok{description}\KeywordTok{:}\AttributeTok{ Analog power supply voltage}
\AttributeTok{    }\FunctionTok{display}\KeywordTok{:}\AttributeTok{ Vdd}
\AttributeTok{    }\FunctionTok{unit}\KeywordTok{:}\AttributeTok{ V}
\AttributeTok{    }\FunctionTok{typical}\KeywordTok{:}\AttributeTok{ }\FloatTok{1.5}
\AttributeTok{  }\FunctionTok{ibias}\KeywordTok{:}
\AttributeTok{    }\FunctionTok{description}\KeywordTok{:}\AttributeTok{ Bias current}
\AttributeTok{    }\FunctionTok{display}\KeywordTok{:}\AttributeTok{ Ibias}
\AttributeTok{    }\FunctionTok{unit}\KeywordTok{:}\AttributeTok{ uA}
\AttributeTok{    }\FunctionTok{typical}\KeywordTok{:}\AttributeTok{ }\DecValTok{20}
\AttributeTok{  }\FunctionTok{corner}\KeywordTok{:}
\AttributeTok{    }\FunctionTok{description}\KeywordTok{:}\AttributeTok{ Process corner}
\AttributeTok{    }\FunctionTok{display}\KeywordTok{:}\AttributeTok{ Corner}
\AttributeTok{    }\FunctionTok{typical}\KeywordTok{:}\AttributeTok{ tt}
\AttributeTok{  }\FunctionTok{temp}\KeywordTok{:}
\AttributeTok{    }\FunctionTok{description}\KeywordTok{:}\AttributeTok{ Ambient temperature}
\AttributeTok{    }\FunctionTok{display}\KeywordTok{:}\AttributeTok{ Temperature}
\AttributeTok{    }\FunctionTok{unit}\KeywordTok{:}\AttributeTok{ °C}
\AttributeTok{    }\FunctionTok{typical}\KeywordTok{:}\AttributeTok{ }\DecValTok{27}
\AttributeTok{  }\FunctionTok{cload}\KeywordTok{:}
\AttributeTok{    }\FunctionTok{description}\KeywordTok{:}\AttributeTok{ Load capacitance}
\AttributeTok{    }\FunctionTok{display}\KeywordTok{:}\AttributeTok{ Cload}
\AttributeTok{    }\FunctionTok{unit}\KeywordTok{:}\AttributeTok{ fF}
\AttributeTok{    }\FunctionTok{typical}\KeywordTok{:}\AttributeTok{ }\DecValTok{50}
\end{Highlighting}
\end{Shaded}

These are the fallback values for every parameter's simulation unless a
specific parameter overrides them. The condition \textbf{name}
(\texttt{vdd}, \texttt{temp}, \texttt{corner}, \ldots) is meaningful: it
must match a \texttt{\$\{vdd\}}-style substitution token used inside
your testbench schematic (see
\protect\hyperlink{4-testbench-templates-and-token-substitution}{§4}).

\hypertarget{parameters-spec-tool-conditions-plot}{%
\subsubsection{3.5 Parameters --- spec, tool, conditions,
plot}\label{parameters-spec-tool-conditions-plot}}

This is where most of the file lives. One example, the AC gain/bandwidth
parameter block:

\begin{Shaded}
\begin{Highlighting}[]
\FunctionTok{parameters}\KeywordTok{:}
\AttributeTok{  }\FunctionTok{ac\_params}\KeywordTok{:}
\AttributeTok{    }\FunctionTok{spec}\KeywordTok{:}
\AttributeTok{      }\FunctionTok{gain}\KeywordTok{:}
\AttributeTok{        }\FunctionTok{display}\KeywordTok{:}\AttributeTok{ Output voltage ratio}
\AttributeTok{        }\FunctionTok{description}\KeywordTok{:}\AttributeTok{ Large{-}signal dc gain}
\AttributeTok{        }\FunctionTok{unit}\KeywordTok{:}\AttributeTok{ V/V}
\AttributeTok{        }\FunctionTok{minimum}\KeywordTok{:}
\AttributeTok{          }\FunctionTok{value}\KeywordTok{:}\AttributeTok{ }\FloatTok{0.97}
\AttributeTok{        }\FunctionTok{typical}\KeywordTok{:}
\AttributeTok{          }\FunctionTok{value}\KeywordTok{:}\AttributeTok{ any}
\AttributeTok{        }\FunctionTok{maximum}\KeywordTok{:}
\AttributeTok{          }\FunctionTok{value}\KeywordTok{:}\AttributeTok{ }\FloatTok{1.03}
\AttributeTok{      }\FunctionTok{bw}\KeywordTok{:}
\AttributeTok{        }\FunctionTok{display}\KeywordTok{:}\AttributeTok{ Bandwidth}
\AttributeTok{        }\FunctionTok{description}\KeywordTok{:}\AttributeTok{ The {-}3dB bandwidth of the buffer}
\AttributeTok{        }\FunctionTok{unit}\KeywordTok{:}\AttributeTok{ Hz}
\AttributeTok{        }\FunctionTok{minimum}\KeywordTok{:}
\AttributeTok{          }\FunctionTok{value}\KeywordTok{:}\AttributeTok{ }\FloatTok{10e6}
\AttributeTok{        }\FunctionTok{typical}\KeywordTok{:}
\AttributeTok{          }\FunctionTok{value}\KeywordTok{:}\AttributeTok{ any}
\AttributeTok{        }\FunctionTok{maximum}\KeywordTok{:}
\AttributeTok{          }\FunctionTok{value}\KeywordTok{:}\AttributeTok{ any}
\AttributeTok{    }\FunctionTok{tool}\KeywordTok{:}
\AttributeTok{      }\FunctionTok{ngspice}\KeywordTok{:}
\AttributeTok{        }\FunctionTok{template}\KeywordTok{:}\AttributeTok{ ota{-}5t{-}ac.sch}
\AttributeTok{        }\FunctionTok{format}\KeywordTok{:}\AttributeTok{ ascii}
\AttributeTok{        }\FunctionTok{suffix}\KeywordTok{:}\AttributeTok{ .data}
\AttributeTok{        }\FunctionTok{variables}\KeywordTok{:}\AttributeTok{ }\KeywordTok{[}\AttributeTok{gain}\KeywordTok{,}\AttributeTok{ bw}\KeywordTok{]}
\AttributeTok{    }\FunctionTok{plot}\KeywordTok{:}
\AttributeTok{      }\FunctionTok{gain\_vs\_corner}\KeywordTok{:}
\AttributeTok{        }\FunctionTok{type}\KeywordTok{:}\AttributeTok{ xyplot}
\AttributeTok{        }\FunctionTok{xaxis}\KeywordTok{:}\AttributeTok{ corner}
\AttributeTok{        }\FunctionTok{yaxis}\KeywordTok{:}\AttributeTok{ gain}
\AttributeTok{        }\FunctionTok{limits}\KeywordTok{:}\AttributeTok{ auto}
\AttributeTok{      }\FunctionTok{bw\_vs\_corner}\KeywordTok{:}
\AttributeTok{        }\FunctionTok{type}\KeywordTok{:}\AttributeTok{ xyplot}
\AttributeTok{        }\FunctionTok{xaxis}\KeywordTok{:}\AttributeTok{ corner}
\AttributeTok{        }\FunctionTok{yaxis}\KeywordTok{:}\AttributeTok{ bw}
\AttributeTok{        }\FunctionTok{limits}\KeywordTok{:}\AttributeTok{ auto}
\AttributeTok{    }\FunctionTok{conditions}\KeywordTok{:}
\AttributeTok{      }\FunctionTok{corner}\KeywordTok{:}
\AttributeTok{        }\FunctionTok{enumerate}\KeywordTok{:}\AttributeTok{ }\KeywordTok{[}\AttributeTok{ss}\KeywordTok{,}\AttributeTok{ sf}\KeywordTok{,}\AttributeTok{ tt}\KeywordTok{,}\AttributeTok{ fs}\KeywordTok{,}\AttributeTok{ ff}\KeywordTok{]}
\AttributeTok{      }\FunctionTok{vdd}\KeywordTok{:}
\AttributeTok{        }\FunctionTok{minimum}\KeywordTok{:}\AttributeTok{ }\FloatTok{1.45}
\AttributeTok{        }\FunctionTok{typical}\KeywordTok{:}\AttributeTok{ }\FloatTok{1.5}
\AttributeTok{        }\FunctionTok{maximum}\KeywordTok{:}\AttributeTok{ }\FloatTok{1.55}
\AttributeTok{      }\FunctionTok{vin}\KeywordTok{:}
\AttributeTok{        }\FunctionTok{minimum}\KeywordTok{:}\AttributeTok{ }\FloatTok{0.7}
\AttributeTok{        }\FunctionTok{typical}\KeywordTok{:}\AttributeTok{ }\FloatTok{0.8}
\AttributeTok{        }\FunctionTok{maximum}\KeywordTok{:}\AttributeTok{ }\FloatTok{0.9}
\AttributeTok{      }\FunctionTok{temp}\KeywordTok{:}
\AttributeTok{        }\FunctionTok{minimum}\KeywordTok{:}\AttributeTok{ }\DecValTok{{-}40}
\AttributeTok{        }\FunctionTok{typical}\KeywordTok{:}\AttributeTok{ }\DecValTok{27}
\AttributeTok{        }\FunctionTok{maximum}\KeywordTok{:}\AttributeTok{ }\DecValTok{130}
\end{Highlighting}
\end{Shaded}

Reading this as a reviewer would:

\begin{itemize}
\tightlist
\item
  \texttt{spec.gain} says: \emph{``I claim 0.97--1.03 V/V unity gain,
  and I will fail my own build if it isn't.''} \texttt{spec.bw} says:
  \emph{``I claim ≥10 MHz bandwidth, no claimed upper bound
  (\texttt{any}).''}
\item
  \texttt{tool.ngspice.template} points at the testbench schematic;
  \texttt{variables} lists which measured quantities to pull out of the
  simulation output.
\item
  \texttt{conditions} is the sweep: five process corners × supply range
  × input range × temperature range, all combined. \texttt{enumerate}
  gives a discrete list (corners);
  \texttt{minimum}/\texttt{typical}/\texttt{maximum} on a numeric
  condition sweeps that variable across its range if a \texttt{step} is
  also given, or just runs the three named points if not.
\item
  \texttt{plot} says which of those sweeps to actually render as a graph
  in the report --- \texttt{gain\_vs\_corner}, \texttt{bw\_vs\_corner},
  etc.
\end{itemize}

\texttt{minimum}/\texttt{typical}/\texttt{maximum} under \texttt{spec}
each accept \texttt{value:\ any} (measure but don't grade), a concrete
\texttt{value} (grade against it, fail if outside),
\texttt{fail:\ false} (measure and report, don't fail the build on it),
and optional \texttt{calculation}/\texttt{limit} overrides for unusual
pass/fail logic (e.g.~``must be above'' vs.~``must be below'').

\begin{center}\rule{0.5\linewidth}{0.5pt}\end{center}

\hypertarget{testbench-templates-and-token-substitution}{%
\subsection{4. Testbench Templates and Token
Substitution}\label{testbench-templates-and-token-substitution}}

Every \texttt{tool.ngspice.template} entry names an xschem schematic
file living under \texttt{cace/templates/}. These are ordinary xschem
schematics, with one twist: anywhere you'd normally write a fixed SPICE
value, you instead write a \textbf{\texttt{CACE\{...\}} token}, which
CACE substitutes at run time with the value for that condition in that
particular sweep point. From the course's actual \texttt{ota-5t-ac.sch}
testbench (inside a \texttt{code\_shown} symbol carrying raw SPICE):

\begin{verbatim}
.include CACE{DUT_path}
.temp CACE{temp}
.param mc_ok = CACE{sigma=1}
.option SEED=CACE[CACE{seed=12345} + CACE{iterations=0}]

.control
set num_threads=1
op
let dcgain=v(v_out)/v(v_in)
ac dec 101 10 100MEG
meas ac acgain MAX vmag(v_out) FROM=10 TO=100
let f3db = acgain/sqrt(2)
meas ac fbw WHEN vmag(v_out)=f3db FALL=1
\end{verbatim}

Three token forms appear here:

\begin{itemize}
\tightlist
\item
  \textbf{\texttt{CACE\{name\}}} --- substitute the current value of
  condition \texttt{name} (e.g.~\texttt{CACE\{temp\}} becomes
  \texttt{27} or \texttt{-40} or \texttt{130} depending on which sweep
  point is running).
\item
  \textbf{\texttt{CACE\{name=default\}}} --- same, but with a fallback
  if the condition isn't defined for this parameter
  (\texttt{CACE\{sigma=1\}} defaults to \texttt{1} --- Monte Carlo
  ``off'' --- unless a Monte Carlo parameter block overrides
  \texttt{sigma}).
\item
  \textbf{\texttt{CACE{[}expression{]}}} --- an arithmetic expression
  combining tokens, evaluated before substitution
  (\texttt{CACE\{seed=12345\}\ +\ CACE\{iterations=0\}} produces a
  distinct SPICE random seed per Monte Carlo iteration).
\end{itemize}

\texttt{DUT\_path} is a CACE-provided pseudo-condition pointing at
whichever generated device-under-test netlist matches the current
\texttt{-s/-\/-source} setting --- schematic, layout, pex, or rcx (see
\protect\hyperlink{8-post-layout-characterization}{§8}). You never
hand-edit that path; it's how the \emph{same} testbench schematic works
unmodified whether you're characterizing the ideal schematic or the
parasitic-extracted layout.

Practical note for the ``software engineer obfuscation'' allergy: there
is no hidden templating engine here beyond straight string substitution
inside a normal SPICE \texttt{.include}/\texttt{.control} block. If you
can read a SPICE deck, you can read a CACE testbench --- the only new
syntax is the \texttt{CACE\{\}}/\texttt{CACE{[}{]}} token itself.

\begin{center}\rule{0.5\linewidth}{0.5pt}\end{center}

\hypertarget{running-cace-and-reading-the-report}{%
\subsection{5. Running CACE and Reading the
Report}\label{running-cace-and-reading-the-report}}

From your project root:

\begin{Shaded}
\begin{Highlighting}[]
\ExtensionTok{cace}\NormalTok{ cace/voltage{-}buffer{-}ota.yaml}
\end{Highlighting}
\end{Shaded}

This runs every parameter block in the datasheet, for every condition
combination, and writes the annotated results into
\texttt{documentation:} (per \protect\hyperlink{32-paths}{§3.2}) as
Markdown, with per-parameter plots as SVG/PNG. Two useful narrowing
flags while iterating:

\begin{Shaded}
\begin{Highlighting}[]
\ExtensionTok{cace}\NormalTok{ cace/voltage{-}buffer{-}ota.yaml {-}p ac\_params        \# only this parameter block}
\ExtensionTok{cace}\NormalTok{ cace/voltage{-}buffer{-}ota.yaml {-}j 4                \# 4 parallel simulation jobs}
\end{Highlighting}
\end{Shaded}

The Markdown report is the artifact you commit to \texttt{docs/}. If you
want an HTML version for easier browsing (not required for submission,
just convenient locally), the course ships a one-line wrapper:

\begin{Shaded}
\begin{Highlighting}[]
\ExtensionTok{cace/cace\_view.sh}\NormalTok{ cace/\_docs/ota{-}5t\_schematic.md}
\end{Highlighting}
\end{Shaded}

--- which is nothing more than \texttt{pandoc\ -f\ markdown\ -t\ html}
followed by opening the result in a browser. No magic there either.

Reading the report: each parameter gets a pass/fail row (measured value
vs.~your \texttt{spec} limits) plus, where a \texttt{plot} block exists,
a graph swept across whichever condition you asked for
(\texttt{gain\_vs\_corner}, \texttt{bw\_vs\_temp}, \ldots). A clean CACE
report with everything green is the single strongest piece of evidence
you can hand a schematic reviewer --- stronger than a paragraph of prose
claiming your circuit works.

\begin{center}\rule{0.5\linewidth}{0.5pt}\end{center}

\hypertarget{pvt-corner-sweeps}{%
\subsection{6. PVT Corner Sweeps}\label{pvt-corner-sweeps}}

``PVT'' = process, voltage, temperature. Doing this by hand --- as your
course material notes --- means either (a) using designer intuition to
guess the worst-case corner and simulating only that, or (b)
brute-forcing every combination with a runner like CACE. Practically
every \texttt{parameters.*.conditions} block you saw in
\protect\hyperlink{35-parameters--spec-tool-conditions-plot}{§3.5} is a
PVT sweep already:

\begin{Shaded}
\begin{Highlighting}[]
\FunctionTok{conditions}\KeywordTok{:}
\AttributeTok{  }\FunctionTok{corner}\KeywordTok{:}
\AttributeTok{    }\FunctionTok{enumerate}\KeywordTok{:}\AttributeTok{ }\KeywordTok{[}\AttributeTok{ss}\KeywordTok{,}\AttributeTok{ sf}\KeywordTok{,}\AttributeTok{ tt}\KeywordTok{,}\AttributeTok{ fs}\KeywordTok{,}\AttributeTok{ ff}\KeywordTok{]}
\AttributeTok{  }\FunctionTok{vdd}\KeywordTok{:}
\AttributeTok{    }\FunctionTok{minimum}\KeywordTok{:}\AttributeTok{ }\FloatTok{1.45}
\AttributeTok{    }\FunctionTok{typical}\KeywordTok{:}\AttributeTok{ }\FloatTok{1.5}
\AttributeTok{    }\FunctionTok{maximum}\KeywordTok{:}\AttributeTok{ }\FloatTok{1.55}
\AttributeTok{  }\FunctionTok{temp}\KeywordTok{:}
\AttributeTok{    }\FunctionTok{minimum}\KeywordTok{:}\AttributeTok{ }\DecValTok{{-}40}
\AttributeTok{    }\FunctionTok{typical}\KeywordTok{:}\AttributeTok{ }\DecValTok{27}
\AttributeTok{    }\FunctionTok{maximum}\KeywordTok{:}\AttributeTok{ }\DecValTok{130}
\end{Highlighting}
\end{Shaded}

Five corners × three supply points × three temperature points = 45
simulations for one parameter block, run automatically and in parallel
(\texttt{-j}). The number of process corners you enumerate is
PDK-dependent --- \texttt{ss/sf/tt/fs/ff} is the standard 5-corner set
(slow-slow, slow-fast, typical-typical, fast-slow, fast-fast) shared
across SG13G2, GF180MCU, and SKY130, though the \emph{library file} you
point at for each differs (see
\protect\hyperlink{9-per-pdk-adaptation-sg13g2cmos5l-gf180mcu-sky130}{§9}).

A callout from your course material worth repeating here, since it
applies directly to a differential buffer/filter block with several
independently-biased stages: \emph{``the number of Monte Carlo
simulation runs \(N\) has to be large enough to approach Gaussian
distributions\ldots{} often \(N=250\) is a good compromise''} --- the
same balancing act applies to how many corner/PVT points you actually
need versus simulation time. Enumerating all 5 corners × full
temperature range for every parameter, every iteration of design, is
often more than you need during early sizing; it becomes necessary once
you're locking down the schematic-review numbers.

\begin{center}\rule{0.5\linewidth}{0.5pt}\end{center}

\hypertarget{monte-carlo-and-mismatch-analysis}{%
\subsection{7. Monte Carlo and Mismatch
Analysis}\label{monte-carlo-and-mismatch-analysis}}

Corners model \emph{global} process variation (an entire wafer skewing
slow or fast together). Monte Carlo with device mismatch models
\emph{local}, random, device-to-device variation --- two
nominally-identical transistors on the same die differing slightly ---
which corners cannot represent and which matters enormously for a
differential input buffer's input-referred offset or a bandgap/Vcm
reference's absolute accuracy.

The course's \texttt{ota-5t.yaml} \texttt{ac\_mc\_params} block is a
complete, working example:

\begin{Shaded}
\begin{Highlighting}[]
\FunctionTok{ac\_mc\_params}\KeywordTok{:}
\AttributeTok{  }\FunctionTok{spec}\KeywordTok{:}
\AttributeTok{    }\FunctionTok{gain\_mc}\KeywordTok{:}
\AttributeTok{      }\FunctionTok{display}\KeywordTok{:}\AttributeTok{ Output voltage ratio (MC)}
\AttributeTok{      }\FunctionTok{unit}\KeywordTok{:}\AttributeTok{ V/V}
\AttributeTok{      }\FunctionTok{minimum}\KeywordTok{:}\AttributeTok{ }\KeywordTok{\{}\AttributeTok{ }\FunctionTok{value}\KeywordTok{:}\AttributeTok{ any }\KeywordTok{\}}
\AttributeTok{      }\FunctionTok{typical}\KeywordTok{:}\AttributeTok{ }\KeywordTok{\{}\AttributeTok{ }\FunctionTok{value}\KeywordTok{:}\AttributeTok{ any }\KeywordTok{\}}
\AttributeTok{      }\FunctionTok{maximum}\KeywordTok{:}\AttributeTok{ }\KeywordTok{\{}\AttributeTok{ }\FunctionTok{value}\KeywordTok{:}\AttributeTok{ any }\KeywordTok{\}}
\AttributeTok{    }\FunctionTok{bw\_mc}\KeywordTok{:}
\AttributeTok{      }\FunctionTok{display}\KeywordTok{:}\AttributeTok{ Bandwidth (MC)}
\AttributeTok{      }\FunctionTok{unit}\KeywordTok{:}\AttributeTok{ Hz}
\AttributeTok{      }\FunctionTok{minimum}\KeywordTok{:}\AttributeTok{ }\KeywordTok{\{}\AttributeTok{ }\FunctionTok{value}\KeywordTok{:}\AttributeTok{ }\FloatTok{10e6}\AttributeTok{ }\KeywordTok{\}}
\AttributeTok{      }\FunctionTok{typical}\KeywordTok{:}\AttributeTok{ }\KeywordTok{\{}\AttributeTok{ }\FunctionTok{value}\KeywordTok{:}\AttributeTok{ any }\KeywordTok{\}}
\AttributeTok{      }\FunctionTok{maximum}\KeywordTok{:}\AttributeTok{ }\KeywordTok{\{}\AttributeTok{ }\FunctionTok{value}\KeywordTok{:}\AttributeTok{ any }\KeywordTok{\}}
\AttributeTok{  }\FunctionTok{tool}\KeywordTok{:}
\AttributeTok{    }\FunctionTok{ngspice}\KeywordTok{:}
\AttributeTok{      }\FunctionTok{template}\KeywordTok{:}\AttributeTok{ ota{-}5t{-}ac.sch}
\AttributeTok{      }\FunctionTok{collate}\KeywordTok{:}\AttributeTok{ iterations}
\AttributeTok{      }\FunctionTok{format}\KeywordTok{:}\AttributeTok{ ascii}
\AttributeTok{      }\FunctionTok{suffix}\KeywordTok{:}\AttributeTok{ .data}
\AttributeTok{      }\FunctionTok{variables}\KeywordTok{:}\AttributeTok{ }\KeywordTok{[}\AttributeTok{gain\_mc}\KeywordTok{,}\AttributeTok{ bw\_mc}\KeywordTok{]}
\AttributeTok{      }\FunctionTok{script}\KeywordTok{:}\AttributeTok{ ota{-}5t{-}ac.py}
\AttributeTok{      }\FunctionTok{script\_variables}\KeywordTok{:}\AttributeTok{ }\KeywordTok{[}\AttributeTok{gain\_mc\_arr}\KeywordTok{,}\AttributeTok{ bw\_mc\_arr}\KeywordTok{]}
\AttributeTok{  }\FunctionTok{plot}\KeywordTok{:}
\AttributeTok{    }\FunctionTok{gain\_mc}\KeywordTok{:}
\AttributeTok{      }\FunctionTok{type}\KeywordTok{:}\AttributeTok{ histogram}
\AttributeTok{      }\FunctionTok{xaxis}\KeywordTok{:}\AttributeTok{ gain\_mc\_arr}
\AttributeTok{    }\FunctionTok{bw\_mc}\KeywordTok{:}
\AttributeTok{      }\FunctionTok{type}\KeywordTok{:}\AttributeTok{ histogram}
\AttributeTok{      }\FunctionTok{xaxis}\KeywordTok{:}\AttributeTok{ bw\_mc\_arr}
\AttributeTok{  }\FunctionTok{conditions}\KeywordTok{:}
\AttributeTok{    }\FunctionTok{iterations}\KeywordTok{:}\CommentTok{               \# for Monte Carlo}
\AttributeTok{      }\FunctionTok{description}\KeywordTok{:}\AttributeTok{ Iterations to run}
\AttributeTok{      }\FunctionTok{minimum}\KeywordTok{:}\AttributeTok{ }\DecValTok{1}
\AttributeTok{      }\FunctionTok{maximum}\KeywordTok{:}\AttributeTok{ }\DecValTok{100}
\AttributeTok{      }\FunctionTok{step}\KeywordTok{:}\AttributeTok{ linear}
\AttributeTok{      }\FunctionTok{stepsize}\KeywordTok{:}\AttributeTok{ }\DecValTok{1}
\AttributeTok{    }\FunctionTok{sigma}\KeywordTok{:}
\AttributeTok{      }\FunctionTok{typical}\KeywordTok{:}\AttributeTok{ }\DecValTok{1}\CommentTok{              \# 1 = mismatch on; 0 = off (see IHP{-}Open{-}PDK issue \#149)}
\AttributeTok{    }\FunctionTok{corner}\KeywordTok{:}
\AttributeTok{      }\FunctionTok{typical}\KeywordTok{:}\AttributeTok{ tt\_stat}\CommentTok{         \# Monte Carlo corner for TT}
\end{Highlighting}
\end{Shaded}

Key mechanics:

\begin{itemize}
\tightlist
\item
  \textbf{\texttt{iterations}} is a swept condition from 1 to 100 (your
  \(N\) from \protect\hyperlink{6-pvt-corner-sweeps}{§6}) --- each
  iteration gets a distinct random seed via the
  \texttt{CACE{[}CACE\{seed=...\}\ +\ CACE\{iterations=...\}{]}}
  expression shown in
  \protect\hyperlink{4-testbench-templates-and-token-substitution}{§4}.
\item
  \textbf{\texttt{collate:\ iterations}} tells CACE to gather all 100
  individual results into arrays (\texttt{gain\_mc\_arr},
  \texttt{bw\_mc\_arr}) rather than reporting 100 separate rows.
\item
  \textbf{\texttt{script:\ ota-5t-ac.py}} is a small user-supplied
  post-processing script (this pattern is documented in CACE's
  \href{https://cace.readthedocs.io/en/latest/tutorials/custom_scripts.html}{``Using
  Custom Scripts to Postprocess Simulation Results''} tutorial) ---
  here, it just collects and possibly computes statistics (mean/σ) from
  the raw per-iteration data before handing it to the histogram plotter.
\item
  \textbf{\texttt{corner:\ tt\_stat}} --- this is the PDK's Monte
  Carlo-enabled variant of the typical corner (naming convention differs
  per PDK, see
  \protect\hyperlink{9-per-pdk-adaptation-sg13g2cmos5l-gf180mcu-sky130}{§9}).
\end{itemize}

A histogram of \texttt{gain\_mc\_arr}/\texttt{bw\_mc\_arr} across
100--250 runs is exactly the evidence a reviewer wants for claims like
``input-referred offset is within ±X mV'' or ``Vcm reference holds to ±Y
mV'' --- the kind of numbers your proof-of-concept Vcm reference and
differential buffers will need. Don't run Monte Carlo and process
corners simultaneously in the same sweep (they answer different
questions --- see
\protect\hyperlink{9-per-pdk-adaptation-sg13g2cmos5l-gf180mcu-sky130}{§9}
for why combining them is usually meaningless); run corners for global
worst-case, Monte Carlo separately for local mismatch/offset statistics.

\begin{center}\rule{0.5\linewidth}{0.5pt}\end{center}

\hypertarget{post-layout-characterization}{%
\subsection{8. Post-Layout
Characterization}\label{post-layout-characterization}}

CACE's \texttt{-s/-\/-source} CLI flag (full reference in
\protect\hyperlink{12-appendix-command-cheat-sheet-and-blank-datasheet-skeleton}{§12})
selects which netlist CACE builds \texttt{DUT\_path} from:

\begin{verbatim}
-s {schematic,layout,pex,rcx,best}
\end{verbatim}

\begin{itemize}
\tightlist
\item
  \texttt{schematic} --- direct from your xschem schematic capture. This
  is what a \textbf{schematic review} wants, and what you should be
  running right now.
\item
  \texttt{layout} --- netlist extracted from layout (via LVS
  extraction), device-level but without parasitics --- mainly a
  connectivity/consistency check against the schematic.
\item
  \texttt{pex} --- layout netlist with \textbf{C}-parasitic extraction
  added (coupling/loading capacitances).
\item
  \texttt{rcx} --- full \textbf{R+C}-parasitic extraction (adds
  parasitic resistance too --- routing IR drop, matching-critical
  resistor mismatch from real geometry).
\item
  \texttt{best} (default) --- use the most complete netlist available,
  falling back to schematic if no layout exists yet.
\end{itemize}

The point that matters for your immediate deadline: \textbf{the exact
same datasheet YAML and testbench templates you write today for the
schematic-review submission are what you'll re-run, unmodified, against
\texttt{pex}/\texttt{rcx} netlists once layout exists.} You are not
writing throwaway schematic-only test infrastructure --- you're writing
the full characterization harness for the entire project, and the
schematic-review milestone is just its first invocation. That's the
strongest argument for investing real effort in the datasheet now rather
than treating it as a review-day formality.

One thing to plan for architecturally, given your proof-of-concept block
mixes fast (gmC oscillator/filter) and slow (bias/Vcm) time constants:
parasitic RC from real layout routing will move pole locations more than
it moves DC operating points, so bandwidth/phase-margin specs are the
ones most likely to visibly shift between \texttt{schematic} and
\texttt{pex}/\texttt{rcx} runs --- worth flagging in your own spec
document (see
\protect\hyperlink{13-the-chipalooza-2-analog-harness-and-what-it-means-for-your-cace-setup}{§13})
as a ``specification difficulty'' rather than discovering it in week 33.

\begin{center}\rule{0.5\linewidth}{0.5pt}\end{center}

\hypertarget{per-pdk-adaptation-sg13g2cmos5l-gf180mcu-sky130}{%
\subsection{9. Per-PDK Adaptation: SG13G2/CMOS5L, GF180MCU,
SKY130}\label{per-pdk-adaptation-sg13g2cmos5l-gf180mcu-sky130}}

Covered in priority order, matching what's time-critical for you right
now.

\hypertarget{ihp-sg13g2-full-stack-and-sg13cmos5l-restricted-primary-immediate}{%
\subsubsection{9.1 IHP SG13G2 (full stack) and SG13CMOS5L (restricted)
--- primary,
immediate}\label{ihp-sg13g2-full-stack-and-sg13cmos5l-restricted-primary-immediate}}

\textbf{The good news first:} SG13CMOS5L is \emph{not} a different
device family. Per the
\href{https://github.com/IHP-GmbH/ihp-sg13cmos5l}{IHP-GmbH/ihp-sg13cmos5l}
repository, it's explicitly the \textbf{M1--M4+TM1 metal stack} (4 thin
routing layers + 1 thick top metal = ``5L'') variant of the same SG13G2
130 nm BiCMOS process your course already uses --- same MOSFETs, same
\texttt{.model} cards, same device symbols in xschem. The metal-stack
reduction is a \textbf{back-end-of-line (layout)} change, not a
schematic/device-model change. That means:

\begin{itemize}
\tightlist
\item
  Your course's existing \texttt{cace/voltage-buffer-ota.yaml},
  testbench templates, and \texttt{xschemrc} setup should port to CMOS5L
  for \textbf{schematic-level} simulation with essentially no changes
  beyond \texttt{PDK:\ ihp-sg13cmos5l} (once that identifier is
  finalized --- see caveat below) instead of \texttt{PDK:\ ihp-sg13g2}.
\item
  The metal-stack restriction only bites once you have layout: fewer
  routing layers, different parasitic RC extraction, and the digital
  flow (LibreLane, per the
  \href{https://github.com/tatzelbrumm/ihp-sg13cmos5l-librelane-template}{\texttt{ihp-sg13cmos5l-librelane-template}})
  needs the reduced stack's LEF/tech files. None of that affects your
  \texttt{\_schematic.md} CACE run.
\end{itemize}

\textbf{The caveat:} the \texttt{ihp-sg13cmos5l} PDK repository's own
README currently states it is \emph{``meant to be used only as a
temporary storage during the development of the build/compile migration
script''} --- i.e.~this is an actively-moving target as of today (22 Aug
2026), not a frozen, versioned PDK release. Practically:

\begin{itemize}
\tightlist
\item
  \textbf{Verify before you rely on it.} Confirm the exact
  \texttt{\$PDK}/\texttt{\$PDK\_ROOT} your Chipalooza toolchain expects,
  and whether MOSFET statistical/mismatch models
  (\texttt{sg13g2\_hv\_nmos\_vfbo\_mm}, \texttt{sg13g2\_hv\_nmos\_rsgo},
  etc. --- the same Gaussian mismatch parameters used in your course's
  \texttt{ota-5t.yaml} \texttt{ac\_mc\_params} block) are confirmed
  present and unchanged in the CMOS5L packaging, since these live in the
  device model files, not the metal stack, but ``should be identical''
  is not the same as ``confirmed identical'' for a PDK still
  mid-migration.
\item
  \textbf{Fallback plan}: if \texttt{ihp-sg13cmos5l} isn't stable enough
  to simulate against directly this week, run your CACE schematic
  characterization against plain \texttt{ihp-sg13g2} (which you already
  have working) and note in your submission that the design targets the
  CMOS5L metal-stack variant of the same process --- schematic-level
  results are expected to be identical since no device models differ.
  Revisit once the PDK repo stabilizes.
\end{itemize}

Corner naming:
\texttt{tt}/\texttt{ss}/\texttt{sf}/\texttt{fs}/\texttt{ff} as in your
course examples; Monte Carlo/mismatch corner is \texttt{tt\_stat} (per
the course's own \texttt{ac\_mc\_params} block, referencing
\href{https://github.com/IHP-GmbH/IHP-Open-PDK/issues/149}{IHP-Open-PDK
issue \#149} for the \texttt{sigma}/\texttt{mc\_ok} semantics).

\hypertarget{gf180mcu-second-priority-time-critical}{%
\subsubsection{9.2 GF180MCU --- second priority,
time-critical}\label{gf180mcu-second-priority-time-critical}}

Not related to CMOS5L by name --- a different foundry family
(GlobalFoundries 180 nm MCU, also offered in 3LM/4LM/5LM metal-stack
variants, which is a separate coincidence of ``5-layer'' terminology
worth not confusing with SG13CMOS5L). If your Chipalooza track (or a
second submission) targets GF180MCU:

\begin{itemize}
\tightlist
\item
  Master device model file: \texttt{sm141064.ngspice} --- BSIM4
  (\texttt{level=54}), covers 3.3 V and 6 V MOSFETs, BJTs, diodes, and
  MIM capacitors.
\item
  \textbf{Corners are selected per device family independently} ---
  separate \texttt{.lib} statements for MOS/resistor/BJT corners,
  e.g.~typical MOS + slow-slow resistor + fast-fast BJT can be mixed in
  one testbench, unlike a single unified corner name.
\item
  \textbf{Monte Carlo/mismatch}: enabled via
  \texttt{.param\ sw\_stat\_mismatch=1} (local device mismatch) and
  optionally \texttt{.param\ sw\_stat\_global=1} (adds inter-die global
  variation on top). A
  \texttt{fet\_mc\_skew}/\texttt{res\_mc\_skew}/\texttt{cap\_mc\_skew}
  parameter (default \texttt{3}) controls how many sigma of
  \emph{global} variation to apply --- lowering it tightens the Monte
  Carlo spread if you have measured data suggesting the default 3σ
  assumption is pessimistic for your process split. Full syntax:
  \href{https://gf180mcu-pdk.readthedocs.io/en/latest/analog/model_parameters/HV/HV_4_1.html}{GF180MCU
  PDK docs, §4.1 Statistical Models}.
\item
  Porting your CACE datasheet: \texttt{default\_conditions.corner}
  enumeration and the testbench's \texttt{.lib} include lines need
  PDK-specific values; the YAML \emph{structure}
  (pins/parameters/spec/conditions) is unchanged.
\end{itemize}

\hypertarget{sky130-third-priority-lower-urgency}{%
\subsubsection{9.3 SKY130 --- third priority, lower
urgency}\label{sky130-third-priority-lower-urgency}}

\begin{itemize}
\tightlist
\item
  Corners selected via
  \texttt{.lib\ "\$PDK\_ROOT/sky130A/libs.tech/ngspice/sky130.lib.spice"\ tt}
  (and \texttt{ff}/\texttt{ss}/\texttt{sf}/\texttt{fs}, plus
  \texttt{leak}/\texttt{wafer} variants for specific analyses).
\item
  \textbf{Monte Carlo/mismatch} uses a \texttt{MC\_MM\_SWITCH} parameter
  (\texttt{.param\ mc\_mm\_switch=1}) wired into \texttt{AGAUSS(...)}
  terms on every BSIM4 card, plus
  \texttt{.option\ SEED=\textless{}value\textgreater{}} for
  reproducible-vs-varying random draws per run --- structurally the same
  idea as SG13G2's \texttt{sigma}/\texttt{mc\_ok} and GF180MCU's
  \texttt{sw\_stat\_mismatch}, different parameter name.
\item
  A useful SKY130-specific shortcut: the pseudo-corner
  \textbf{\texttt{tt\_mm}} is ``typical corner + local mismatch
  enabled'' in one \texttt{.lib} selection, rather than needing a
  separate switch --- check whether an equivalent shortcut corner exists
  for SG13G2/CMOS5L or GF180MCU before assuming you need the two-step
  (\texttt{corner=tt} + separate mismatch switch) pattern everywhere.
\item
  Explicit warning worth restating for all three PDKs, not just SKY130:
  \textbf{don't run process corners and Monte Carlo mismatch in the same
  sweep.} Corners already assume every device shifts together in the
  worst direction; stacking random per-device mismatch on top of an
  already-worst-case corner represents a vanishingly low-probability
  combined event, not a realistic one, and will make your spec look
  artificially bad (or you'll be tempted to loosen limits to compensate,
  which defeats the purpose). Run them as separate \texttt{parameters}
  blocks, as the course's own \texttt{ac\_params} (corners)
  vs.~\texttt{ac\_mc\_params} (Monte Carlo, fixed at \texttt{tt\_stat})
  split demonstrates.
\end{itemize}

\hypertarget{summary-table}{%
\subsubsection{9.4 Summary table}\label{summary-table}}

\begin{longtable}[]{@{}llll@{}}
\toprule
\begin{minipage}[b]{0.22\columnwidth}\raggedright
\strut
\end{minipage} & \begin{minipage}[b]{0.22\columnwidth}\raggedright
SG13G2 / SG13CMOS5L\strut
\end{minipage} & \begin{minipage}[b]{0.22\columnwidth}\raggedright
GF180MCU\strut
\end{minipage} & \begin{minipage}[b]{0.22\columnwidth}\raggedright
SKY130\strut
\end{minipage}\tabularnewline
\midrule
\endhead
\begin{minipage}[t]{0.22\columnwidth}\raggedright
Device models\strut
\end{minipage} & \begin{minipage}[t]{0.22\columnwidth}\raggedright
Same BSIM-type models for both stack variants\strut
\end{minipage} & \begin{minipage}[t]{0.22\columnwidth}\raggedright
\texttt{sm141064.ngspice}, BSIM4 (level 54)\strut
\end{minipage} & \begin{minipage}[t]{0.22\columnwidth}\raggedright
BSIM4 (level 54), per-corner \texttt{.lib} files\strut
\end{minipage}\tabularnewline
\begin{minipage}[t]{0.22\columnwidth}\raggedright
Corner selection\strut
\end{minipage} & \begin{minipage}[t]{0.22\columnwidth}\raggedright
Single \texttt{.lib}/corner name
(\texttt{tt}/\texttt{ss}/\texttt{sf}/\texttt{fs}/\texttt{ff})\strut
\end{minipage} & \begin{minipage}[t]{0.22\columnwidth}\raggedright
Independent per device family (MOS/R/BJT)\strut
\end{minipage} & \begin{minipage}[t]{0.22\columnwidth}\raggedright
Single \texttt{.lib} file, corner as argument\strut
\end{minipage}\tabularnewline
\begin{minipage}[t]{0.22\columnwidth}\raggedright
Mismatch enable\strut
\end{minipage} & \begin{minipage}[t]{0.22\columnwidth}\raggedright
\texttt{sigma}/\texttt{mc\_ok} param (course convention)\strut
\end{minipage} & \begin{minipage}[t]{0.22\columnwidth}\raggedright
\texttt{sw\_stat\_mismatch=1} (+ \texttt{sw\_stat\_global=1})\strut
\end{minipage} & \begin{minipage}[t]{0.22\columnwidth}\raggedright
\texttt{mc\_mm\_switch=1} + \texttt{AGAUSS()}\strut
\end{minipage}\tabularnewline
\begin{minipage}[t]{0.22\columnwidth}\raggedright
MC ``typical'' shortcut\strut
\end{minipage} & \begin{minipage}[t]{0.22\columnwidth}\raggedright
\texttt{tt\_stat} corner\strut
\end{minipage} & \begin{minipage}[t]{0.22\columnwidth}\raggedright
--- (compose manually)\strut
\end{minipage} & \begin{minipage}[t]{0.22\columnwidth}\raggedright
\texttt{tt\_mm} corner\strut
\end{minipage}\tabularnewline
\begin{minipage}[t]{0.22\columnwidth}\raggedright
Metal stack caveat\strut
\end{minipage} & \begin{minipage}[t]{0.22\columnwidth}\raggedright
CMOS5L repo is WIP as of Aug 2026 --- verify before relying on it\strut
\end{minipage} & \begin{minipage}[t]{0.22\columnwidth}\raggedright
3LM/4LM/5LM variants exist; confirm which your shuttle uses\strut
\end{minipage} & \begin{minipage}[t]{0.22\columnwidth}\raggedright
Single standard stack\strut
\end{minipage}\tabularnewline
\bottomrule
\end{longtable}

\begin{center}\rule{0.5\linewidth}{0.5pt}\end{center}

\hypertarget{model-accuracy-caveats}{%
\subsection{10. Model-Accuracy Caveats}\label{model-accuracy-caveats}}

Worth stating plainly rather than assuming: \textbf{open-source PDK
Monte Carlo and mismatch models are, in general, less mature and less
silicon-validated than a commercial foundry's qualified design kit.}
They are typically derived from a mix of foundry-supplied statistical
corners, published process variation data, and community
reverse-engineering/curve-fitting --- not decades of proprietary
characterization across millions of production die. This isn't a reason
to distrust them wholesale, but it is a reason to treat MC/mismatch
\emph{numbers} with more margin than you would in a commercial-PDK
tapeout, and to say so explicitly in your submission rather than
implicitly overclaiming precision the model can't actually back up.

Practical mitigations, roughly in order of effort:

\begin{enumerate}
\def\labelenumi{\arabic{enumi}.}
\tightlist
\item
  \textbf{Add explicit margin} on any spec that leans on a Monte
  Carlo/mismatch result (offset, PSRR, absolute Vcm accuracy) rather
  than taking the simulated 3σ number at face value as your final
  claimed spec.
\item
  \textbf{Cross-check against published silicon data where it exists}
  --- for SG13G2 in particular, IHP publishes process specification
  documents (e.g.~the
  \href{https://github.com/IHP-GmbH/IHP-Open-PDK/blob/main/ihp-sg13g2/libs.doc/doc/SG13G2_os_process_spec.pdf}{\texttt{SG13G2\_os\_process\_spec.pdf}})
  with measured parameter spreads; sanity-check your Monte Carlo
  histogram's σ against those where the parameter overlaps
  (e.g.~threshold voltage mismatch).
\item
  \textbf{Use layout structures matched to known-good references where
  possible} --- your own proposal already leans this direction (``derive
  the pmos and nmos bias currents from layout structures that match the
  layout of the current and voltage references on the Chipalooza
  harness''), which sidesteps absolute-model-accuracy risk by relying on
  \emph{matching} (a much better-modeled effect than absolute value)
  rather than absolute device parameters.
\item
  \textbf{Flag known-uncertain parameters explicitly in your submission}
  rather than silently shipping a tight spec --- a reviewer trusts ``we
  simulate ±5 mV offset with a known ±2× model-confidence caveat, so
  we're specifying ±15 mV'' far more than an unqualified ``±5 mV'' that
  turns out to be a simulation artifact.
\end{enumerate}

\begin{center}\rule{0.5\linewidth}{0.5pt}\end{center}

\hypertarget{mapping-cace-output-to-chipalooza-schematic-review-expectations}{%
\subsection{11. Mapping CACE Output to Chipalooza Schematic-Review
Expectations}\label{mapping-cace-output-to-chipalooza-schematic-review-expectations}}

Putting §§1--10 together into what actually goes in front of a reviewer:

\begin{itemize}
\tightlist
\item
  \textbf{The schematic}
  (\texttt{xschem/\textless{}name\textgreater{}.sch},
  \texttt{\textless{}name\textgreater{}\_tb.sch}), matching the pin list
  in your CACE datasheet exactly (CACE will already have caught any
  mismatch before you got this far, since it checks pin names against
  the schematic).
\item
  \textbf{The datasheet}
  (\texttt{cace/\textless{}name\textgreater{}.yaml}), which \emph{is}
  your written specification --- a reviewer should be able to read
  \texttt{spec:} blocks and know exactly what you're claiming without
  reading a separate prose document.
\item
  \textbf{The generated report}
  (\texttt{docs/\textless{}name\textgreater{}\_schematic.md}), run with
  \texttt{-s\ schematic}, showing every claimed parameter passing (or
  explicitly flagged/explained if not) across the PVT and Monte Carlo
  sweeps from §§6--7.
\item
  \textbf{A short human-readable description} of what the block does,
  its architecture, and any known specification difficulties or open
  questions (your proposal document is already exactly this format) ---
  CACE doesn't replace this, it substantiates it.
\end{itemize}

A reviewer's real question at this stage is not ``does this look like a
plausible circuit'' (anyone can draw a plausible-looking OTA) but ``has
this designer actually verified, under a defined and reproducible set of
conditions, that the circuit does what they say it does.'' A clean,
committed CACE run answers that question in a way prose alone cannot.

\begin{center}\rule{0.5\linewidth}{0.5pt}\end{center}

\hypertarget{appendix-command-cheat-sheet-and-blank-datasheet-skeleton}{%
\subsection{12. Appendix: Command Cheat-Sheet and Blank Datasheet
Skeleton}\label{appendix-command-cheat-sheet-and-blank-datasheet-skeleton}}

\hypertarget{commands}{%
\subsubsection{12.1 Commands}\label{commands}}

\begin{Shaded}
\begin{Highlighting}[]
\CommentTok{\# Run everything in the datasheet}
\ExtensionTok{cace}\NormalTok{ cace/}\OperatorTok{\textless{}}\NormalTok{name}\OperatorTok{\textgreater{}}\NormalTok{.yaml}

\CommentTok{\# Run one parameter block only, useful while iterating}
\ExtensionTok{cace}\NormalTok{ cace/}\OperatorTok{\textless{}}\NormalTok{name}\OperatorTok{\textgreater{}}\NormalTok{.yaml {-}p ac\_params}

\CommentTok{\# Explicit netlist source (schematic review =\textgreater{} schematic)}
\ExtensionTok{cace}\NormalTok{ cace/}\OperatorTok{\textless{}}\NormalTok{name}\OperatorTok{\textgreater{}}\NormalTok{.yaml {-}s schematic}

\CommentTok{\# Parallel jobs}
\ExtensionTok{cace}\NormalTok{ cace/}\OperatorTok{\textless{}}\NormalTok{name}\OperatorTok{\textgreater{}}\NormalTok{.yaml {-}j 4}

\CommentTok{\# Convert schematic to a standalone SPICE netlist (already done internally by cace,}
\CommentTok{\# but useful to inspect what it generated)}
\CommentTok{\# {-}\textgreater{} see cace/netlist/schematic/\textless{}name\textgreater{}.spice after a run}

\CommentTok{\# Convert a CACE Markdown report to HTML for local viewing (course helper script)}
\ExtensionTok{cace/cace\_view.sh}\NormalTok{ cace/\_docs/}\OperatorTok{\textless{}}\NormalTok{name}\OperatorTok{\textgreater{}}\NormalTok{\_schematic.md}
\end{Highlighting}
\end{Shaded}

Full CLI reference: \texttt{cace\ -\/-help}, or
\href{https://cace.readthedocs.io/en/latest/usage/cace_cli.html}{cace.readthedocs.io/en/latest/usage/cace\_cli.html}.

\hypertarget{blank-datasheet-skeleton}{%
\subsubsection{12.2 Blank datasheet
skeleton}\label{blank-datasheet-skeleton}}

\begin{Shaded}
\begin{Highlighting}[]
\FunctionTok{name}\KeywordTok{:}\AttributeTok{           \textless{}cell\_name\textgreater{}}
\FunctionTok{description}\KeywordTok{:}\AttributeTok{    \textless{}one{-}line description\textgreater{}}
\FunctionTok{PDK}\KeywordTok{:}\AttributeTok{            \textless{}pdk\_identifier\textgreater{}}
\FunctionTok{cace\_format}\KeywordTok{:}\AttributeTok{    }\FloatTok{5.2}

\FunctionTok{authorship}\KeywordTok{:}
\AttributeTok{  }\FunctionTok{designer}\KeywordTok{:}\AttributeTok{         \textless{}your name\textgreater{}}
\AttributeTok{  }\FunctionTok{email}\KeywordTok{:}\AttributeTok{            \textless{}your email\textgreater{}}
\AttributeTok{  }\FunctionTok{creation\_date}\KeywordTok{:}\AttributeTok{    \textless{}date\textgreater{}}
\AttributeTok{  }\FunctionTok{license}\KeywordTok{:}\AttributeTok{          Apache 2.0}

\FunctionTok{paths}\KeywordTok{:}
\AttributeTok{  }\FunctionTok{root}\KeywordTok{:}\AttributeTok{             ..}
\AttributeTok{  }\FunctionTok{schematic}\KeywordTok{:}\AttributeTok{        xschem}
\AttributeTok{  }\FunctionTok{netlist}\KeywordTok{:}\AttributeTok{          cace/netlist}
\AttributeTok{  }\FunctionTok{documentation}\KeywordTok{:}\AttributeTok{    cace/\_docs}
\AttributeTok{  }\FunctionTok{runs}\KeywordTok{:}\AttributeTok{             \_runs}

\FunctionTok{pins}\KeywordTok{:}
\AttributeTok{  }\FunctionTok{\textless{}pin\_name\textgreater{}}\KeywordTok{:}
\AttributeTok{    }\FunctionTok{description}\KeywordTok{:}\AttributeTok{ \textless{}text\textgreater{}}
\AttributeTok{    }\FunctionTok{type}\KeywordTok{:}\AttributeTok{ power|ground|signal|digital}
\AttributeTok{    }\FunctionTok{direction}\KeywordTok{:}\AttributeTok{ input|output|inout}
\CommentTok{    \# Vmin / Vmax / Imin / Imax as needed}

\FunctionTok{default\_conditions}\KeywordTok{:}
\AttributeTok{  }\FunctionTok{\textless{}condition\_name\textgreater{}}\KeywordTok{:}
\AttributeTok{    }\FunctionTok{description}\KeywordTok{:}\AttributeTok{ \textless{}text\textgreater{}}
\AttributeTok{    }\FunctionTok{display}\KeywordTok{:}\AttributeTok{ \textless{}short label\textgreater{}}
\AttributeTok{    }\FunctionTok{unit}\KeywordTok{:}\AttributeTok{ \textless{}unit\textgreater{}}
\AttributeTok{    }\FunctionTok{typical}\KeywordTok{:}\AttributeTok{ \textless{}value\textgreater{}}

\FunctionTok{parameters}\KeywordTok{:}
\AttributeTok{  }\FunctionTok{\textless{}parameter\_block\_name\textgreater{}}\KeywordTok{:}
\AttributeTok{    }\FunctionTok{spec}\KeywordTok{:}
\AttributeTok{      }\FunctionTok{\textless{}result\_name\textgreater{}}\KeywordTok{:}
\AttributeTok{        }\FunctionTok{display}\KeywordTok{:}\AttributeTok{ \textless{}label\textgreater{}}
\AttributeTok{        }\FunctionTok{description}\KeywordTok{:}\AttributeTok{ \textless{}text\textgreater{}}
\AttributeTok{        }\FunctionTok{unit}\KeywordTok{:}\AttributeTok{ \textless{}unit\textgreater{}}
\AttributeTok{        }\FunctionTok{minimum}\KeywordTok{:}\AttributeTok{ }\KeywordTok{\{}\AttributeTok{ }\FunctionTok{value}\KeywordTok{:}\AttributeTok{ \textless{}value\textgreater{}|any }\KeywordTok{\}}
\AttributeTok{        }\FunctionTok{typical}\KeywordTok{:}\AttributeTok{  }\KeywordTok{\{}\AttributeTok{ }\FunctionTok{value}\KeywordTok{:}\AttributeTok{ \textless{}value\textgreater{}|any }\KeywordTok{\}}
\AttributeTok{        }\FunctionTok{maximum}\KeywordTok{:}\AttributeTok{ }\KeywordTok{\{}\AttributeTok{ }\FunctionTok{value}\KeywordTok{:}\AttributeTok{ \textless{}value\textgreater{}|any }\KeywordTok{\}}
\AttributeTok{    }\FunctionTok{tool}\KeywordTok{:}
\AttributeTok{      }\FunctionTok{ngspice}\KeywordTok{:}
\AttributeTok{        }\FunctionTok{template}\KeywordTok{:}\AttributeTok{ \textless{}testbench\textgreater{}.sch}
\AttributeTok{        }\FunctionTok{format}\KeywordTok{:}\AttributeTok{ ascii}
\AttributeTok{        }\FunctionTok{suffix}\KeywordTok{:}\AttributeTok{ .data}
\AttributeTok{        }\FunctionTok{variables}\KeywordTok{:}\AttributeTok{ }\KeywordTok{[}\AttributeTok{\textless{}result\_name\textgreater{}}\KeywordTok{,}\AttributeTok{ ...}\KeywordTok{]}
\AttributeTok{    }\FunctionTok{plot}\KeywordTok{:}
\AttributeTok{      }\FunctionTok{\textless{}plot\_name\textgreater{}}\KeywordTok{:}
\AttributeTok{        }\FunctionTok{type}\KeywordTok{:}\AttributeTok{ xyplot|histogram|semilogx|semilogy|loglog}
\AttributeTok{        }\FunctionTok{xaxis}\KeywordTok{:}\AttributeTok{ \textless{}condition\_or\_variable\textgreater{}}
\AttributeTok{        }\FunctionTok{yaxis}\KeywordTok{:}\AttributeTok{ \textless{}result\_name\textgreater{}}
\AttributeTok{        }\FunctionTok{limits}\KeywordTok{:}\AttributeTok{ auto}
\AttributeTok{    }\FunctionTok{conditions}\KeywordTok{:}
\AttributeTok{      }\FunctionTok{corner}\KeywordTok{:}
\AttributeTok{        }\FunctionTok{enumerate}\KeywordTok{:}\AttributeTok{ }\KeywordTok{[}\AttributeTok{ss}\KeywordTok{,}\AttributeTok{ sf}\KeywordTok{,}\AttributeTok{ tt}\KeywordTok{,}\AttributeTok{ fs}\KeywordTok{,}\AttributeTok{ ff}\KeywordTok{]}
\CommentTok{      \# other swept conditions here}
\end{Highlighting}
\end{Shaded}

\begin{center}\rule{0.5\linewidth}{0.5pt}\end{center}

\hypertarget{the-chipalooza-2-analog-harness-and-what-it-means-for-your-cace-setup}{%
\subsection{13. The Chipalooza \#2 Analog Harness, and What It Means for
Your CACE
Setup}\label{the-chipalooza-2-analog-harness-and-what-it-means-for-your-cace-setup}}

Your project folder already contains the harness reference
(\texttt{sg13cmos5l\_chipalooza\_harness.pdf}) --- the actual
infrastructure your \texttt{s2d\_d2s\_pinbuffers} block plugs into.
Reading it against your proposal's pin table changes a few ``to be
discussed'' items from open questions into constraints you can design
against, and it changes how you should structure the CACE
\texttt{pins:}/\texttt{default\_conditions:} sections.

\hypertarget{what-the-harness-actually-provides}{%
\subsubsection{13.1 What the harness actually
provides}\label{what-the-harness-actually-provides}}

\begin{itemize}
\tightlist
\item
  \textbf{Per-slot dedicated analog I/O}: two pins,
  \texttt{sN\_an{[}1:0{]}}, per project slot --- matches your proposal's
  \texttt{vin}/\texttt{vout} dedicated-pin concept directly (one slot,
  two dedicated analog pins).
\item
  \textbf{Shared analog bus}: \texttt{analog{[}2:0{]}} --- three shared
  analog lines, routed to any slot through an analog switch matrix
  (\texttt{ana\_route\_sel}). This is the harness-level realization of
  your proposal's \texttt{vcm}/\texttt{vdiffp}/\texttt{vdiffn} ``shared
  analog pins'' concept --- they are not extra pins your block owns,
  they're a shared resource your block's switch logic taps into.
\item
  \textbf{Bias/Vcm generation is centralized, not something your block
  must synthesize from scratch}: each slot gets
  \texttt{sN\_ibias{[}1{]}}, \texttt{sN\_ibias{[}2{]}}, and
  \texttt{sN\_vbias}, sourced from two shared iDACs and a shared voltage
  bias generator, configured over SPI (registers
  \texttt{0x15}--\texttt{0x1A}). Your proposal's ``to be discussed,
  100nA ballpark'' \texttt{ibias} and ``to be discussed, if adjustable''
  \texttt{vbias} are, per the harness, \textbf{digitally trimmable from
  a shared generator} --- which changes the design question from ``how
  do I generate a stable 100 nA on-chip'' to ``what range and resolution
  do I need the shared iDAC to give me, and how sensitive is my circuit
  to its trim step size.''
\item
  \textbf{Power domains match your table almost exactly}: independently
  gated \texttt{vdd3v3}/\texttt{vss3v3} (analog) and
  \texttt{vdd1v2}/\texttt{vss1v2} (digital, can run at 1.5 V) per slot,
  separate from the I/O ring's \texttt{vddio}/\texttt{vssio} ---
  confirms your proposal's pin table rather than requiring changes.
\item
  \textbf{Digital control goes through a shared crossbar, not dedicated
  pins}: \texttt{gpio\_di{[}23:0{]}} in / \texttt{gpio\_do{[}7:0{]}}
  out, assigned per-slot via SPI registers (\texttt{0x20}--\texttt{0x37}
  input assignments, \texttt{0x40}--\texttt{0x4f} output assignments)
  rather than your scan chain owning fixed physical pins. Your
  \texttt{ena}/\texttt{reset}/\texttt{outbuf\_en}/\texttt{inbuf\_en}/\texttt{filter\_en}/\texttt{vdiff\_en{[}1:0{]}}/\texttt{vcmsel{[}1:0{]}}/\texttt{scanclk}/\texttt{scanctrl{[}1:0{]}}/\texttt{scanin}/\texttt{scanout}
  signals are all internal-to-your-block-boundary signals that get
  \emph{bound} to specific \texttt{gpio\_di}/\texttt{gpio\_do} bit
  positions by harness configuration, not fixed pins your schematic
  hardwires to package pins.
\end{itemize}

\hypertarget{on-analog20-vs.-a-possible-analog30}{%
\subsubsection{\texorpdfstring{13.2 On \texttt{analog{[}2:0{]}} vs.~a
possible
\texttt{analog{[}3:0{]}}}{13.2 On analog{[}2:0{]} vs.~a possible analog{[}3:0{]}}}\label{on-analog20-vs.-a-possible-analog30}}

Per your instruction: treat a 4th shared analog line
(\texttt{analog{[}3{]}}) as \emph{possible} --- the harness may be
revised before final submission --- but don't build a dependency on it
existing. Concretely, in your CACE setup:

\begin{itemize}
\tightlist
\item
  \textbf{Don't hardcode ``3 shared analog pins'' as a structural
  assumption anywhere load-bearing.} Your \texttt{pins:} block should
  describe your own block's boundary (however many
  \texttt{vdiffp}/\texttt{vdiffn}/\texttt{vcm}-equivalent pins
  \emph{your circuit} exposes) --- the harness's bus width is a
  routing-matrix property external to your block, not a parameter your
  datasheet needs to know at all. As long as your block's differential
  I/O count doesn't itself change, \texttt{analog{[}2:0{]}} becoming
  \texttt{analog{[}3:0{]}} is invisible to your CACE datasheet and
  testbenches.
\item
  \textbf{Where it \emph{does} matter} is your standalone/breakout test
  plan (your proposal's final section) --- if you're relying on shared
  analog lines to route unbuffered internal differential signals out for
  external measurement (as you describe, for measuring VCM and the
  unbuffered filter output), a wider shared bus gives you more
  simultaneous internal test-node visibility. Write that section so it
  degrades gracefully with 3 lines (your 9 listed tests already
  time-multiplex the same 2--3 shared lines across different
  configurations) rather than assuming a 4th line will be available to
  run two of those tests simultaneously.
\item
  \textbf{Practical hedge}: if you do get a firm answer before final
  submission, updating from 3 to 4 shared lines is a
  harness-side/testbench-side change only (which physical
  \texttt{analog{[}n{]}} your test setup taps), not a change to your
  circuit's own \texttt{pins:} list or its CACE parameter specs.
\end{itemize}

\hypertarget{consequence-for-your-datasheets-default_conditions}{%
\subsubsection{\texorpdfstring{13.3 Consequence for your datasheet's
\texttt{default\_conditions}}{13.3 Consequence for your datasheet's default\_conditions}}\label{consequence-for-your-datasheets-default_conditions}}

Because \texttt{ibias}/\texttt{vbias} are harness-supplied and digitally
trimmed rather than internally generated, model them in your CACE
datasheet the same way the course's \texttt{ota-5t.yaml} models
\texttt{ibias\_20u} --- as an \textbf{input pin with a condition sweep},
not as an internal node you're trying to characterize the accuracy of.
Your own circuit's job is to be correctly biased \emph{given} whatever
the harness iDAC delivers (with its own trim resolution/range as a swept
condition,
e.g.~\texttt{ibias:\ \{\ minimum:\ 80,\ typical:\ 100,\ maximum:\ 120,\ unit:\ nA\ \}}
reflecting expected trim tolerance), not to prove the iDAC itself is
accurate --- that's the harness team's characterization burden, not
yours. This is a meaningful scope reduction for your ``specification
difficulties'' section: it turns ``generate a stable low-noise 100 nA
reference'' from your problem into ``operate correctly across the iDAC's
specified trim range,'' which is a much smaller and more tractable CACE
sweep.

\begin{center}\rule{0.5\linewidth}{0.5pt}\end{center}

\hypertarget{sources}{%
\subsection{Sources}\label{sources}}

\begin{itemize}
\tightlist
\item
  \href{https://github.com/efabless/cace}{efabless/cace} --- CACE tool
  repository
\item
  \href{https://cace.readthedocs.io/en/latest/}{CACE Documentation} ---
  overview, tutorials, reference manual
\item
  \href{https://cace.readthedocs.io/en/latest/reference/datasheet_format.html}{CACE
  Datasheet Format reference}
\item
  \href{https://cace.readthedocs.io/en/latest/usage/cace_cli.html}{CACE
  Command Line Interface reference}
\item
  \href{https://cace.readthedocs.io/en/latest/tutorials/custom_scripts.html}{CACE
  custom-scripts tutorial}
\item
  \href{https://github.com/efabless/sky130_ef_ip__template}{efabless/sky130\_ef\_ip\_\_template}
  --- analog IP repository convention
\item
  \href{https://github.com/efabless/sky130_sw_ip__por}{efabless/sky130\_sw\_ip\_\_por}
  --- real Chipalooza IP example (bandgap PoR)
\item
  \href{https://github.com/IHP-GmbH/ihp-sg13cmos5l}{IHP-GmbH/ihp-sg13cmos5l}
  --- SG13CMOS5L PDK repository (WIP status)
\item
  \href{https://github.com/tatzelbrumm/ihp-sg13cmos5l-librelane-template}{tatzelbrumm/ihp-sg13cmos5l-librelane-template}
  --- LibreLane digital-flow template for SG13CMOS5L
\item
  \href{https://github.com/IHP-GmbH/IHP-Open-PDK/issues/149}{IHP-GmbH/IHP-Open-PDK
  issue \#149} --- \texttt{sigma}/\texttt{mc\_ok} Monte Carlo semantics
\item
  \href{https://github.com/IHP-GmbH/IHP-Open-PDK/blob/main/ihp-sg13g2/libs.doc/doc/SG13G2_os_process_spec.pdf}{IHP-Open-PDK:
  SG13G2\_os\_process\_spec.pdf}
\item
  \href{https://gf180mcu-pdk.readthedocs.io/en/latest/analog/model_parameters/HV/HV_4_1.html}{GF180MCU
  PDK docs --- Statistical Models Syntax \& Usage}
\item
  H. Pretl, M. Koefinger, S. Dorrer, \emph{Analog (Integrated) Circuit
  Design}, Johannes Kepler University (Apache-2.0), and its
  \texttt{cace/voltage-buffer-ota.yaml},
  \texttt{cace/templates/ota-5t-ac.sch}, \texttt{cace/cace\_view.sh} ---
  the working CACE example this walkthrough is built around
\item
  \texttt{sg13cmos5l\_chipalooza\_harness.pdf} (user-provided) ---
  Chipalooza \#2 analog harness specification
\end{itemize}

\end{document}
